\section{Weight--monodromy, Hyodo--Kato cohomology, and rigid analytic
formality}\label{weightmonodromy-hyodokato-cohomology-and-rigid-analytic-formality}

\textbf{Scope and provenance.} Literature cutoff: 20 November 2026.
Checked on 21 November 2026 UTC by consequences\_criteria, GPT-6 Astra
Pro; reasoning effort is not exposed. The adopted RH and Dirichlet-GRH
inputs are not needed for these local geometric results. The sources
below are additional geometric inputs. The 2026 preprints are used as
cited theorems, with their statements and indicated dependency arguments
inspected; their complete foundational proofs have not been
independently reconstructed.

\subsection{CC-WM01 --- Exact known classes of
weight--monodromy}\label{cc-wm01-exact-known-classes-of-weightmonodromy}

Fix the convention that geometric Frobenius \(F\) and monodromy satisfy
\(FNF^{-1}=q^{-1}N\), where \(q\) is the residue-field cardinality.
Monodromy purity of weight \(i\) means that eigenvalues on
\(\operatorname{gr}^M_j\) are algebraic and have absolute value
\(q^{(i+j)/2}\) under every complex embedding. The monodromy filtration
is centered at zero.

\begin{longtable}[]{@{}
  >{\raggedright\arraybackslash}p{(\columnwidth - 4\tabcolsep) * \real{0.3333}}
  >{\raggedright\arraybackslash}p{(\columnwidth - 4\tabcolsep) * \real{0.3333}}
  >{\raggedright\arraybackslash}p{(\columnwidth - 4\tabcolsep) * \real{0.3333}}@{}}
\toprule\noalign{}
\begin{minipage}[b]{\linewidth}\raggedright
Primary result and inspected locator
\end{minipage} & \begin{minipage}[b]{\linewidth}\raggedright
Exact relevant hypotheses
\end{minipage} & \begin{minipage}[b]{\linewidth}\raggedright
Conclusion and boundary
\end{minipage} \\
\midrule\noalign{}
\endhead
\bottomrule\noalign{}
\endlastfoot
Ito and Zou (2026), arXiv:2610.07890v1, \textbf{Theorem B}, p.~2;
\textbf{Theorem 4.4}, pp.~12--13 & \(K/\mathbf Q_p\) finite; \(X/K\)
hyper-Kähler: smooth projective, geometrically simply connected, with
\(H^0(X,\Omega^2_{X/K})=K\sigma\) for a nowhere-degenerate two-form;
definition, p.~1 & Weight--monodromy for every \(0\le i\le2\dim X\) and
every prime \(\ell\), including \(p\), via the associated Weil--Deligne
representation. \textbf{No \(b_2\ge4\) hypothesis occurs in this
theorem}; that restriction concerns separate compatibility results. \\
Scholze (2012), arXiv:1111.4914v1, \textbf{Theorem 9.6}, p.~49;
Conjecture 9.3, p.~48 & \(k\) a characteristic-zero local field; \(Y/k\)
geometrically connected, proper and smooth; \(Y\) a
\textbf{set-theoretic} complete intersection in a projective smooth
toric variety & Weight--monodromy in all degrees for \(\ell\ne p\). \\
Binda, Kato, and Vezzani (2025), arXiv:2207.00369v2, \textbf{Theorem
1.2}, p.~2; \textbf{Theorem 5.6}, p.~28 & \(K/\mathbf Q_p\) finite;
\(Y/K\) smooth and a \textbf{scheme-theoretic} complete intersection in
a projective smooth toric variety & The p-adic conjecture in every
degree. Theorem 5.6 more generally treats such a variety over
\(\mathbf C_p\), asserting quasi-purity of
\(H^i_{\mathrm{HK}}(Y^{\mathrm{an}})\). The
set-theoretic/scheme-theoretic distinction is explicit on p.~2. \\
Wear (2023), arXiv:2303.05610v1, \textbf{Theorem 1.1.2}, p.~3;
coefficient convention, pp.~1--2 & \(k\) a local field of residue
characteristic \(p\); \(A/k\) abelian; \(Y\) geometrically connected,
proper and smooth, and a \textbf{set-theoretic} complete intersection in
\(A\) & Weight--monodromy in all degrees for \(\ell\ne p\). This permits
equal or mixed characteristic and does not provide the p-adic case. \\
\end{longtable}

Scholze's \textbf{Theorem 9.4}, p.~48, states Deligne's
equal-characteristic result for \(k\) the local field at
\(x\in C(\mathbf F_q)\), where \(C/\mathbf F_q\) is a curve and the
variety comes from a proper smooth scheme over \(C\setminus\{x\}\). This
exactly scoped globally defined case suffices in his reduction.

The inspected Ito--Zou dependency chain is \textbf{Lemma 3.8}, p.~8
(purity in degrees at most two), \textbf{Lemma 4.2}, p.~12 (the
higher-degree tensor square is a summand of tensors of \(H^2\) and its
dual), and \textbf{Lemma 4.3}, p.~12 (tensor-square detection), followed
by Theorem 4.4. The complete last argument was read; its earlier
semisimplicity and Mumford--Tate inputs remain cited results.

In a semistable Hyodo--Kato model with \(q=p^a\), use the linear
operator \(\Phi^a\), rather than semilinear \(\Phi\), for traces. The
relation \(N\Phi=p\Phi N\) yields \(\Phi^aN(\Phi^a)^{-1}=q^{-1}N\).
Binda--Kato--Vezzani's p.~2 comparison discussion identifies the
pertinent potentially semistable p-adic realization; their Theorem 5.6's
final splitting argument was read, retaining the preceding motivic
comparison theorems as inputs.

\subsection{CC-WM02 --- The cochain-algebra
consequence}\label{cc-wm02-the-cochain-algebra-consequence}

Petrov and Zavyalov (2026), arXiv:2607.14517v2, \textbf{Theorem 1.2},
p.~2, is \textbf{Theorem 5.1 and Lemma 5.2}, p.~18. For \(X/K\) smooth
proper rigid analytic, \(K/\mathbf Q_p\) finite, weight--monodromy in
all degrees for a fixed prime \(\ell\), possibly \(p\), implies

\[
 R\Gamma_{\mathrm{et}}(X_{\mathbf C_p},\mathbf Q_\ell)
 \simeq \bigoplus_i H^i_{\mathrm{et}}(X_{\mathbf C_p},\mathbf Q_\ell)[-i]
\]

as underlying \(E_\infty\)-algebras. For \(\ell=p\), de Rham cochains
are formal too. Definitions 5.5--5.6, pp.~19--20, specify monodromy
purity; Examples 5.4, p.~19, use \(D_{\mathrm{pst}}\) at \(p\). No
canonical or Galois-equivariant splitting is asserted. The final proof
is on pp.~25--26.

Smooth proper \textbf{algebraic} varieties already have formal cochains
unconditionally, as the source explains on p.~2. Thus the substantive
additional application is to genuinely rigid analytic spaces. Its
\textbf{Construction 3.1 and Proposition 3.2}, p.~5, provide a smooth
proper rigid surface with nonformal cochains at every prime and in de
Rham cohomology. Smoothness and properness alone therefore cannot give
the needed hypothesis.

\subsection{CC-WM03 --- A finite family of exterior trace conditions
implies rigid
formality}\label{cc-wm03-a-finite-family-of-exterior-trace-conditions-implies-rigid-formality}

\textbf{Proposition.} Let \(K/\mathbf Q_p\) be finite and \(X/K\)
smooth, proper and rigid analytic. Fix a prime \(\ell\), including the
possibility \(\ell=p\). For each nonzero cohomology group let

\[
 V_i=D\bigl(H^i_{\mathrm{et}}(X_{\mathbf C_p},\mathbf Q_\ell)\bigr),
 \qquad d_i=\dim V_i,
\]

be the Weil--Deligne representation in Petrov--Zavyalov's Examples 5.4.
Choose a finite extension \(L/K\) killing the finite inertial actions in
these finitely many representations. Put \(q=\#k_L\), and write
\((F_i,N_i)\) for their Frobenius and monodromy over \(L\). Extend
scalars to an algebraic closure of the coefficient field and fix an
embedding of \(\overline{\mathbf Q}\) into that closure. Suppose that
all \(F_i\)-eigenvalues lie in its image, so they can be regarded as
algebraic numbers. For \(1\le m\le d_i\), define

\[
 T_{i,m}(r)=\operatorname{Tr}\left(F_i^r\mid
       \ker N_{\wedge^m V_i}\right),\qquad
 T^\vee_{i,m}(r)=\operatorname{Tr}\left((F_i^\vee)^r\mid
       \ker N_{\wedge^m V_i^\vee}\right).
\]

Here \(F_i^\vee\) denotes inverse-transpose action, and \(N\) acts as a
derivation on exterior powers, with \(N^\vee=-{}^tN\) on the dual. For
every \(i,m\) and embedding
\(\sigma:\overline{\mathbf Q}\hookrightarrow\mathbf C\), suppose that
positive constants \(C_{i,m,\sigma},C^\vee_{i,m,\sigma}\) and
nonnegative exponents \(A_{i,m,\sigma},A^\vee_{i,m,\sigma}\) exist such
that, for \textbf{every integer \(r\ge1\)},

\[
 \begin{split}
 |\sigma(T_{i,m}(r))|&\le C_{i,m,\sigma}r^{A_{i,m,\sigma}}q^{rmi/2},\\
 |\sigma(T^\vee_{i,m}(r))|&\le C^\vee_{i,m,\sigma}r^{A^\vee_{i,m,\sigma}}q^{-rmi/2}.
 \end{split} \tag{CC-WM.1}
\]

Then weight--monodromy holds in every degree over \(K\), and the
underlying \(E_\infty\)-algebra
\(R\Gamma_{\mathrm{et}}(X_{\mathbf C_p},\mathbf Q_\ell)\) is formal. For
\(\ell=p\), \(R\Gamma_{\mathrm{dR}}(X/K)\) is formal as well.
Conversely, weight--monodromy implies algebraicity and (CC-WM.1), with
exponents zero and constants equal to the respective kernel dimensions
when these dimensions are positive.

\textbf{Complete proof of the deduction.} There are finitely many
nonzero finite-dimensional cohomology groups. Intersect their open
inertia kernels and pass to a finite extension to obtain \(L\). Each
pair satisfies \(F_iN_iF_i^{-1}=q^{-1}N_i\).

Fix \(i\). Replacing \(F_i\) by its semisimple part leaves all relevant
traces and monodromy-graded eigenvalues unchanged, by the first part of
BR-WM-02. The same proof decomposes the pair into invariant Jordan
strings, with bases \(e_0,\ldots,e_{s-1}\) such that \(Ne_j=e_{j+1}\),
\(Ne_{s-1}=0\), and \(Fe_j=\mu q^{-j}e_j\). Since every \(\mu\) is
algebraic, these matrices define a pair over a number field. Thus every
\(\sigma\) can be applied to the pair, without requiring an algebraic
model for the full original Weil action. Exterior powers, duals and
kernels commute with this characteristic-zero field extension; their
power traces are precisely the algebraic numbers in (CC-WM.1).

Apply BR-WM-02 to this complex pair with dimension \(d_i\) and weight
\(w=i\). Its trace criterion gives

\[
 |\sigma(\alpha)|=q^{(i+j)/2}
 \quad\text{for every eigenvalue }\alpha
 \text{ on }\operatorname{gr}^M_jV_i.
\]

Doing this for every \(\sigma\) proves monodromy purity over \(L\),
including algebraicity. This is the hypothesis in Petrov--Zavyalov's
Definition 5.6.

For completeness, purity descends from \(L\) to \(K\). Write
\(q_0=\#k_K\). For any chosen geometric Frobenius \(\phi\) over \(K\),
some positive power \(\phi^f\) belongs to the open subgroup \(W_L\);
replacing \(L\) by a further finite extension if necessary does not
affect purity, since its Frobenius is a positive power after the finite
inertial action has been killed. The element \(\phi^f\) has degree \(f\)
over \(K\), so its eigenvalues have absolute value \(q_0^{f(i+j)/2}\) on
\(\operatorname{gr}^M_j\). Indeed, over \(L\) the finite inertial action
is trivial and its degree is \(f/[k_L:k_K]\); hence it acts by that
power of an \(L\)-Frobenius. If \(\alpha\) is a \(\phi\)-eigenvalue,
\(\alpha^f\) is such an eigenvalue. Algebraicity descends from
\(\alpha^f\) by the equation \(t^f-\alpha^f=0\), and every complex
absolute value of \(\alpha\) is \(q_0^{(i+j)/2}\). Scaling \(N\) under a
change of monodromy normalization does not change its filtration. This
proves descent; ascent follows by the same power calculation after
killing finite inertia.

Petrov--Zavyalov's Theorem 5.1 now gives étale formality. At \(p\),
Lemma 5.2 gives de Rham formality. The \(p\)-adic pair is obtained from
\(D_{\mathrm{pst}}\); when a semistable Hyodo--Kato description is
available, use its linear Frobenius and \(N\). No equality with ordinary
\(p\)-adic Galois inertia invariants is needed.

Conversely, monodromy purity gives algebraicity and BR-WM-02's first
condition under every complex embedding. That theorem bounds eigenvalues
on the two invariant exterior spaces by \(q^{mi/2}\) and \(q^{-mi/2}\).
A power trace is the sum of the eigenvalue powers, with algebraic
multiplicities, also for a nonsemisimple Frobenius. The triangle
inequality proves (CC-WM.1) with zero exponents and the stated
constants. A zero cohomology group requires no conditions; a zero kernel
admits any positive constant. Every assertion follows. \(\square\)

\textbf{Exact gain and limit.} There are at most \(2\sum_i d_i\) tested
representations; the quantifier over Frobenius powers remains infinite.
This is a criterion and its formality consequence, not a proof of the
bounds for arbitrary \(X\). For \(\ell\ne p\), after unipotent Galois
inertia is achieved, the kernels are actual inertia invariants of the
exterior representations. At \(p\), they belong to the Weil--Deligne or
Hyodo--Kato realization. Alternating traces on total cohomology do not
supply the separate inequalities. Both dual bounds and all complex
embeddings are required. The adopted Dirichlet GRH does not, by itself,
establish these geometric bounds.

\subsection{Exact primary bibliography and version
record}\label{exact-primary-bibliography-and-version-record}

\begin{itemize}
\tightlist
\item
  Binda, F., Kato, H., \& Vezzani, A. (2025). On the \(p\)-adic
  weight-monodromy conjecture for complete intersections in toric
  varieties. \emph{Inventiones Mathematicae, 241}, 559--603.
  https://doi.org/10.1007/s00222-025-01344-x. Inspected version:
  https://arxiv.org/pdf/2207.00369v2, dated 10 June 2025, 36 pages;
  locators above refer to this version. Published online 17 June 2025.
\item
  Ito, K., \& Zou, H. (2026). \emph{Semisimplicity, purity and
  Mumford-Tate conjecture for hyper-Kähler varieties}
  (arXiv:2610.07890v1) {[}Preprint{]}.
  https://arxiv.org/abs/2610.07890v1. Submitted 6 October 2026, 07:38:56
  UTC. Inspected PDF: https://arxiv.org/pdf/2610.07890v1, 23 pages.
\item
  Petrov, A., \& Zavyalov, B. (2026). \emph{Formality for rigid-analytic
  spaces satisfying the weight-monodromy conjecture}
  (arXiv:2607.14517v2) {[}Preprint{]}.
  https://arxiv.org/abs/2607.14517v2. First submitted 16 July 2026; v2
  revised 17 August 2026, 21:05:01 UTC. Inspected PDF:
  https://arxiv.org/pdf/2607.14517v2, 28 pages.
\item
  Scholze, P. (2012). Perfectoid spaces. \emph{Publications
  Mathématiques de l'IHÉS, 116}, 245--313.
  https://doi.org/10.1007/s10240-012-0042-x. Inspected author version:
  https://arxiv.org/pdf/1111.4914v1, submitted 21 November 2011,
  16:14:23 UTC, 51 pages; numbered locators above refer to that author
  version.
\item
  Wear, P. (2023). \emph{Perfectoid covers of abelian varieties and the
  weight-monodromy conjecture} (arXiv:2303.05610v1) {[}Preprint{]}.
  https://arxiv.org/abs/2303.05610v1. Submitted 9 March 2023, 22:41:32
  UTC. Inspected PDF: https://arxiv.org/pdf/2303.05610v1, 54 pages.
\end{itemize}
